In this lecture, we turn to distributed computing. We define the synchronous complete network model,
study the leader election problem, and analyze a randomized algorithm that elects a leader in $\bigO{1}$
rounds with only $\bigO{\sqrt{n} \log^{\sfrac{3}{2}} n}$ messages. We then take a first look at random
walks on undirected graphs, and we prove the weak law of large numbers with Chebyshev's inequality.

\section{Distributed Computing in a Complete Network}

The network is a \emph{complete graph} $G$ on $n$ nodes (\cref{fig:network-0908}). Every pair of nodes
is joined by a bidirectional communication link, so there are $\binom{n}{2} = \frac{n(n-1)}{2}$ edges
and every node has $n - 1$ incident edges. The network is \emph{synchronous}: all nodes share a clock,
and the computation proceeds in \emph{rounds}. In one round, every node receives the messages sent
to it in the previous round, does a local computation, and sends messages on some of its edges
(receives $\to$ computes $\to$ sends).

\begin{figure}[ht]
    \centering
    \begin{tikzpicture}[scale=0.82]
        % complete network
        \node[point, label=above left:{\small $1$}] (a1) at (-1.3,0.9) {};
        \node[point, label=above:{\small $2$}] (a2) at (0,1.8) {};
        \node[point, label=above right:{\small $3$}] (a3) at (1.3,0.9) {};
        \node[point, label=below right:{\small $n-1$}] (a4) at (0.8,-0.6) {};
        \node[point, label=below left:{\small $n$}] (a5) at (-0.8,-0.6) {};
        \foreach \x/\y in {a1/a2, a1/a3, a1/a4, a1/a5, a2/a3, a2/a4, a2/a5, a3/a4, a3/a5, a4/a5} {
                \draw (\x) -- (\y);
            }
        % candidates and referees
        \begin{scope}[xshift=6cm]
            \node[point, label=left:{\small $u$}] (u) at (0,1.3) {};
            \node[point, label=left:{\small $v$}] (v) at (0,-0.1) {};
            \foreach \i/\h in {1/1.9, 2/1.25, 3/0.6, 4/-0.05, 5/-0.7} {
                    \node[point] (r\i) at (3.2,\h) {};
                }
            \node[right] at (3.3,0.6) {\small $r$};
            \foreach \i in {1,2} {\draw[->, shorten >=2pt] (u) -- (r\i);}
            \foreach \i in {4,5} {\draw[->, shorten >=2pt] (v) -- (r\i);}
            % the two arrows into the common referee end slightly apart
            \draw[->, shorten >=2pt] (u) -- ([yshift=3pt]r3.west);
            \draw[->, shorten >=2pt] (v) -- ([yshift=-3pt]r3.west);
            \node at (0,2.4) {\small candidates};
            \node at (3.2,2.4) {\small referees};
        \end{scope}
    \end{tikzpicture}
    \caption{Left: the complete network on $n$ nodes. Right: candidates $u$ and $v$ with a common
        referee $r$.}
    \label{fig:network-0908}
\end{figure}

\subsection{Leader Election}

\begin{problem}[Leader Election]
    All nodes run the same algorithm. At the end, exactly one node must be
    in the state \emph{leader}, and all other nodes must be in the state \emph{non-leader}.
\end{problem}

The goals are to minimize the number of rounds and the total number of messages sent. The nodes are
either \emph{labelled}, i.e., every node has a unique id, or \emph{anonymous}, i.e., the nodes have
no ids and cannot be told apart.

\paragraph{The trivial algorithm.} With unique ids, every node sends its id to all other nodes in the
first round, and the node with the minimum id becomes the leader. This takes $\bigO{1}$ rounds and
$n(n-1) = \bigTh{n^2}$ messages.

\paragraph{Anonymous nodes.} A node addresses its edges by $n - 1$ \emph{ports}, numbered
$1, \dots, n-1$, one per incident edge. It can send a message on a given port and sees on which port
a message arrives, but it does not know which node is at the other end of a port, and the port
numbering is chosen by an adversary.

\needspace{9\baselineskip} % keep the exercise and the start of its solution on one page
\begin{exercise}\label{ex:anonymous-0908}
    Prove that there is no deterministic algorithm that elects a leader among anonymous nodes in a
    complete network, no matter how many rounds the algorithm is allowed to run.
\end{exercise}

The solution is \cref{thm:anonymous-0908}.

\begin{theorem}[No Deterministic Leader Election for Anonymous Nodes]\label{thm:anonymous-0908}
    Let $n \geq 2$. In an anonymous complete network in which all nodes start in the same state, there
    is no deterministic algorithm that elects a leader for every port numbering: for every
    deterministic algorithm there is a port numbering for which, after every number of rounds, the
    number of leaders is either $0$ or $n$.
\end{theorem}

\begin{proof}
    Fix a deterministic algorithm. Identify the nodes with $\mathbb{Z}_n = \set{0, 1, \dots, n-1}$
    (addition modulo $n$), and let port $j$ of node $i$ lead to node $i + j$, for $j = 1, \dots, n-1$.
    This is a valid port numbering, since $i + 1, \dots, i + (n-1)$ are exactly the nodes other than
    $i$. The edge between $i$ and $i + j$ is port $j$ at $i$ and port $n - j$ at $i + j$, since
    $(i + j) + (n - j) = i$ in $\mathbb{Z}_n$, so a message sent on port $j$ arrives on port $n - j$.

    We show by induction on $t$ that at the end of round $t$ all nodes are in the same state, and
    that in round $t$ all nodes send the same message $m_j$ on port $j$, for every $j$ (``no message''
    counts as a message). For $t = 0$ all nodes are in the same initial state, and nothing has been
    sent. If the claim holds for round $t - 1$, then in round $t$ node $i$ receives on port $q$ the
    message that node $i + q$ sent on its port $n - q$, which is the common message $m_{n-q}$ of
    round $t - 1$ and does not depend on $i$. So all nodes are in the same state and receive the
    same messages on the same ports, and since the algorithm is deterministic, they compute the same
    new state and send the same messages on the same ports.

    Whether a node declares itself the leader is a function of its state. Therefore, at the end of
    every round, either all $n \geq 2$ nodes are leaders or no node is, and there is never exactly
    one leader.
\end{proof}

From now on we assume that the nodes have unique ids.

\paragraph{Complexity classes.} Deterministic, bounded-error probabilistic and bounded-error quantum
polynomial time satisfy $\mathrm{P} \subseteq \mathrm{BPP} \subseteq \mathrm{BQP}$.

\subsection{Known Bounds}

For leader election in a complete network of $n$ nodes with unique ids, the following bounds are
known ($\widetilde{O}$ and $\widetilde{\Omega}$ hide factors polynomial in $\log n$). We prove only
the first and the third row.
\begin{center}
    \begin{tabular}{lll}
        Messages                                           & Rounds            & Algorithm                \\
        \hline
        $\bigTh{n^2}$                                      & $\bigO{1}$        & trivial, deterministic   \\
        $\bigO{n \log n}$                                  & $\bigO{\log n}$   & deterministic            \\
        $\bigO{\sqrt{n} \log^{\sfrac{3}{2}} n} = \widetilde{O}(\sqrt{n})$ & $\bigO{1}$ & randomized (this lecture) \\
        $\widetilde{\Omega}(\sqrt{n})$                     &                   & classical lower bound    \\
        $\widetilde{O}(n^{\sfrac{1}{3}})$                  &                   & quantum                  \\
    \end{tabular}
\end{center}

\section{A Randomized Algorithm for Leader Election}

In this section, $\log$ is the natural logarithm, $p = \frac{2 \log n}{n}$ and
$k = \left\lceil \sqrt{n \log n} \right\rceil$. We have $p \leq 1$ and $k \leq n$ for every
$n \geq 1$, since $2 \log n \leq n$. The algorithm has four steps.
\begin{enumerate}
    \item \textbf{Subsample candidates.} Every node becomes a \emph{candidate} with probability $p$,
          independently of all other nodes.
    \item \textbf{Candidates contact referees.} Every candidate sends its id to $k$ random nodes,
          its \emph{referees}.
    \item \textbf{Referee round.} Every node that received at least one id declares the minimum id
          among the received messages the \emph{winner}, and replies with the winner's id to all
          candidates that contacted it.
    \item \textbf{Decision.} A candidate becomes the leader if it is the winner at all of its
          referees. All other nodes are non-leaders.
\end{enumerate}

\paragraph{The model for the random choices.} Every node $u$ draws a coin $C_u \in \set{0,1}$ with
$\prob{C_u = 1} = p$ ($u$ is a candidate if $C_u = 1$) and a \emph{referee set} $R_u$, which is a
uniformly random set of $k$ distinct nodes among all $n$ nodes. All the $C_u$ and all the $R_u$ are
mutually independent. Only candidates use their referee sets. A candidate may pick itself as a
referee; then it plays the role of this referee itself, and no message is needed.

\subsection{The Number of Candidates}

Let $X = \sum_{u} C_u$ be the number of candidates. By linearity of expectation,
$\expectation{X} = np = 2 \log n$.

\begin{exercise}\label{ex:atleast-0908}
    Show that, with high probability, there is at least $1$ leader.
\end{exercise}

The solution is \cref{thm:leader-atleast-0908}.

\begin{theorem}[At Least One Leader]\label{thm:leader-atleast-0908}
    There is at least one leader if and only if there is at least one candidate, and
    \[\prob{\text{no leader}} = \prob{X = 0} \leq \frac{1}{n^2},
        \qquad\text{i.e.,}\qquad
        \prob{\text{at least $1$ leader}} \geq 1 - \frac{1}{n^2}.\]
\end{theorem}

\begin{proof}
    Only candidates become leaders, so if $X = 0$ there is no leader. Conversely, if $X \geq 1$, let
    $w$ be the candidate with the minimum id (unique, since the ids are unique). Every referee $r$
    of $w$ receives only ids of candidates, all at least the id of $w$, and the id of $w$ is one of
    them, so $r$ declares $w$ the winner. Hence $w$ is a leader. The coins are independent, so,
    using $0 \leq 1 - p$ and $1 - x \leq e^{-x}$,
    $\prob{X = 0} = (1 - p)^{n} \leq e^{-pn} = e^{-2 \log n} = \sfrac{1}{n^2}$.
\end{proof}

\subsection{At Most One Leader}

\begin{exercise}\label{ex:atmost-0908}
    Show that, with probability tending to $1$, there is at most $1$ leader.
\end{exercise}

The solution is \cref{thm:leader-atmost-0908}, which uses \cref{lem:badpair-0908,lem:disjoint-0908};
together with \cref{thm:leader-atleast-0908} it gives exactly one leader with probability tending
to $1$ (\cref{thm:leader-0908}). Call a pair of different nodes $\set{u,v}$ a
\emph{bad pair} if both are candidates and they have no common referee, i.e., $C_u = C_v = 1$ and
$R_u \cap R_v = \emptyset$.

\begin{lemma}[Two Leaders Form a Bad Pair]\label{lem:badpair-0908}
    If two candidates have a common referee, then at most one of them is a leader. Consequently, if
    there is more than one leader, then there is a bad pair.
\end{lemma}

\begin{proof}
    Let $u \neq v$ be candidates with a common referee $r$, where the id of $u$ is smaller than the
    id of $v$ (\cref{fig:network-0908}). The referee $r$ receives both ids (if $r$ is $u$ or $v$
    itself, it knows its own id), so the minimum id that it receives is at most the id of $u$. Hence
    $r$ does not declare $v$ the winner, and $v$ is not a leader. So two different leaders are two
    candidates without a common referee, i.e., a bad pair.
\end{proof}

\begin{lemma}[Two Referee Sets Rarely Miss Each Other]\label{lem:disjoint-0908}
    For any two different nodes $u$ and $v$,
    \[\prob{R_u \cap R_v = \emptyset} \leq \left(1 - \frac{k}{n}\right)^{k}
        \leq e^{-\sfrac{k^2}{n}} \leq \frac{1}{n}.\]
\end{lemma}

\begin{proof}
    Condition on $R_u = S$, where $S$ is any set of $k$ nodes; by independence, $R_v$ is still a
    uniformly random set of $k$ nodes. If $2k > n$, two sets of $k$ nodes cannot be disjoint and the
    probability is $0$. So let $2k \leq n$, and draw the $k$ nodes of $R_v$ one after the other, each
    uniformly among the nodes not drawn yet. If the first $i$ nodes drawn are outside $S$, then
    $n - i$ nodes are left and $k$ of them are in $S$, so the next node is outside $S$ with
    probability $\frac{n - k - i}{n - i}$. Hence
    \[\prob{R_v \cap S = \emptyset}
        = \prod_{i=0}^{k-1} \left(1 - \frac{k}{n - i}\right)
        \leq \left(1 - \frac{k}{n}\right)^{k}
        \leq e^{-\sfrac{k^2}{n}} \leq e^{-\log n} = \frac{1}{n},\]
    because every factor satisfies $0 \leq 1 - \frac{k}{n-i} \leq 1 - \frac{k}{n}$ (as $i \leq k - 1$
    and $2k \leq n$ give $n - i \geq k$), and by $1 - x \leq e^{-x}$ and $k^2 \geq n \log n$. The
    bound does not depend on $S$, so it also holds without conditioning.
\end{proof}

\begin{theorem}[At Most One Leader]\label{thm:leader-atmost-0908}
    \[\prob{\text{more than $1$ leader}} \leq \frac{2 \log^2 n}{n}.\]
\end{theorem}

\begin{proof}
    Let $B$ be the number of bad pairs. For a fixed pair $\set{u,v}$, the coins $C_u$, $C_v$ and the
    pair of sets $(R_u, R_v)$ are independent, so by \cref{lem:disjoint-0908} the pair is bad with
    probability $p^2 \cdot \prob{R_u \cap R_v = \emptyset} \leq \sfrac{p^2}{n}$. By linearity of
    expectation over the $\binom{n}{2}$ pairs,
    $\expectation{B} \leq \binom{n}{2} \frac{p^2}{n} = \frac{n(n-1)}{2} \cdot \frac{4 \log^2 n}{n^3}
        \leq \frac{2 \log^2 n}{n}$.
    If there is more than one leader, then $B \geq 1$ by \cref{lem:badpair-0908}. Since $B \geq 0$,
    Markov's inequality (\cref{thm:markov}) with $b = 1$ gives
    $\prob{\text{more than $1$ leader}} \leq \prob{B \geq 1} \leq \expectation{B}$.
\end{proof}

\begin{theorem}[Exactly One Leader]\label{thm:leader-0908}
    The algorithm elects exactly one leader with probability at least
    \[1 - \frac{1}{n^2} - \frac{2 \log^2 n}{n}.\]
\end{theorem}

\begin{proof}
    The algorithm fails if there is no leader or more than one leader. By the union bound and
    \cref{thm:leader-atleast-0908,thm:leader-atmost-0908}, this has probability at most
    $\frac{1}{n^2} + \frac{2 \log^2 n}{n}$.
\end{proof}

\subsection{Messages and Rounds}

\begin{theorem}[Complexity of the Algorithm]\label{thm:complexity-0908}
    The algorithm takes $3$ rounds, and messages are sent only in the first two. The expected number
    of messages is at most $4 k \log n = \bigO{\sqrt{n} \log^{\sfrac{3}{2}} n}$.
\end{theorem}

\begin{proof}
    \emph{Rounds.} In round $1$ every node draws its coin and its referee set, and the candidates
    send their ids. In round $2$ the referees receive the ids and reply with the minimum. In round
    $3$ the candidates receive the replies and decide; no message is sent.

    \emph{Messages.} Every candidate sends at most $k$ messages and gets at most $k$ replies, so the
    number of messages is at most $2kX$, with expectation at most $2k \expectation{X} = 4 k \log n$,
    which is $\bigO{\sqrt{n} \log^{\sfrac{3}{2}} n}$ since $k \leq \sqrt{n \log n} + 1$.
\end{proof}

\section{Markov Chains: Random Walks in Undirected Graphs}

\begin{definition}[Random Walk]
    Let $G$ be an undirected graph in which every node has at least one neighbour. A \textbf{random
    walk} on $G$ is the following process: a walker starts at some node, and in every step it moves
    from its current node $u$ to a neighbouring node chosen with uniform probability, i.e., to each
    neighbour of $u$ with probability $\frac{1}{\deg(u)}$, independently of the earlier steps.
\end{definition}

We describe the walk by the row vector $p_t = (p_t^1, \dots, p_t^n)$, where $p_t^i$ is the probability
that the walker is at node $i$ after $t$ steps. The walker is at one node, so the entries
of $p_t$ sum to $1$. If the walker is at node $j$ after $t$ steps, it moves to a given neighbour $i$
of $j$ with probability $\frac{1}{\deg(j)}$, so by the law of total probability
\begin{equation}\label{eq:walkstep-0908}
    p_{t+1}^i = \sum_{j \,:\, j \text{ is a neighbour of } i} \frac{p_t^j}{\deg(j)}.
\end{equation}

\begin{figure}[ht]
    \centering
    \begin{tikzpicture}[scale=0.82]
        % six-node example
        \node[point, label=below left:{\small $1$}] (b1) at (0,0) {};
        \node[point, label=left:{\small $2$}] (b2) at (-0.2,1.6) {};
        \node[point, label=above:{\small $3$}] (b3) at (1.4,2.4) {};
        \node[point, label=right:{\small $4$}] (b4) at (3.3,1.7) {};
        \node[point, label=below right:{\small $5$}] (b5) at (2.7,0) {};
        \node[point, label=right:{\small $6$}] (b6) at (2.25,1.05) {};
        \foreach \x/\y in {b1/b2, b1/b3, b1/b5, b1/b6, b2/b3, b3/b4, b3/b6, b4/b5, b4/b6, b5/b6} {
                \draw (\x) -- (\y);
            }
        % triangle
        \begin{scope}[xshift=7cm, yshift=0.3cm]
            \node[point, label=left:{\small $1$}] (c1) at (0,0) {};
            \node[point, label=above:{\small $2$}] (c2) at (1.2,1.8) {};
            \node[point, label=right:{\small $3$}] (c3) at (2.4,0) {};
            \draw (c1) -- (c2) -- (c3) -- (c1);
        \end{scope}
    \end{tikzpicture}
    \caption{Left: a graph with $6$ nodes and $10$ edges. Right: the triangle.}
    \label{fig:walks-0908}
\end{figure}

\begin{example}
    In the graph on the left of \cref{fig:walks-0908}, a walker that starts at node $1$ has
    $p_0 = (1, 0, 0, 0, 0, 0)$ and, since node $1$ has the $4$ neighbours $2$, $3$, $5$, $6$,
    $p_1 = \left(0, \tfrac{1}{4}, \tfrac{1}{4}, 0, \tfrac{1}{4}, \tfrac{1}{4}\right)$. In the triangle
    on the right every node has degree $2$, and starting at node~$1$, by \cref{eq:walkstep-0908},
    $p_0 = (1, 0, 0)$, $p_1 = \left(0, \tfrac{1}{2}, \tfrac{1}{2}\right)$ and
    $p_2 = \left(\tfrac{1}{2}, \tfrac{1}{4}, \tfrac{1}{4}\right)$.
\end{example}

\begin{exercise}\label{ex:triangle-0908}
    Compute $p_t = \left(p_t^1, p_t^2, p_t^3\right)$ for all $t \geq 0$ for the random walk on the
    triangle that starts at node $1$.
\end{exercise}

The solution is \cref{lem:triangle-0908}.

\begin{lemma}[Random Walk on the Triangle]\label{lem:triangle-0908}
    For the random walk on the triangle that starts at node $1$, for all $t \geq 0$,
    \[p_t^1 = \frac{1}{3} + \frac{2}{3} \left(-\frac{1}{2}\right)^{t},
        \qquad
        p_t^2 = p_t^3 = \frac{1}{3} - \frac{1}{3} \left(-\frac{1}{2}\right)^{t}.\]
    In particular, $p_t$ converges to $\left(\frac{1}{3}, \frac{1}{3}, \frac{1}{3}\right)$ as $t \to \infty$.
\end{lemma}

\begin{proof}
    Let $i$ be any of the three nodes, and let $j$ and $l$ be the other two. Both are neighbours of
    $i$ and have degree $2$, so by \cref{eq:walkstep-0908}, and since $p_t^i + p_t^j + p_t^l = 1$,
    $p_{t+1}^i = \frac{1}{2} \left(p_t^j + p_t^l\right) = \frac{1}{2} \left(1 - p_t^i\right)$, and hence
    $p_{t+1}^i - \frac{1}{3} = -\frac{1}{2} \left(p_t^i - \frac{1}{3}\right)$.
    By induction on $t$, $p_t^i = \frac{1}{3} + \left(p_0^i - \frac{1}{3}\right) \left(-\frac{1}{2}\right)^{t}$.
    With $p_0^1 = 1$ and $p_0^2 = p_0^3 = 0$ this gives the two formulas, and
    $\left(-\frac{1}{2}\right)^{t} \to 0$ gives the limit.
\end{proof}

\section{The Law of Large Numbers}

\begin{exercise}\label{ex:wlln-0908}
    Show the ``law of large numbers'' with Chebyshev's inequality.
\end{exercise}

The solution is \cref{thm:wlln-0908}, the \emph{weak} law of large numbers.

\begin{theorem}[Weak Law of Large Numbers]\label{thm:wlln-0908}
    Let $X_1, X_2, \dots$ be pairwise independent random variables with the same expectation
    $\expectation{X_i} = \mu$ and the same finite variance $\variance{X_i} = \sigma^2$ (for example,
    independent and identically distributed with finite variance), and let
    $\overline{X}_n = \frac{1}{n} \sum_{i=1}^{n} X_i$. Then, for every $\eps > 0$,
    \[\prob{\abs{\overline{X}_n - \mu} \geq \eps} \leq \frac{\sigma^2}{n \eps^2}
        \longrightarrow 0 \qquad \text{as } n \to \infty.\]
\end{theorem}

\begin{proof}
    By linearity of expectation, $\expectation{\overline{X}_n} = \frac{1}{n} \sum_{i} \expectation{X_i} = \mu$.
    Write $Y_i = X_i - \mu$, so that $\expectation{Y_i} = 0$, $\expectation{Y_i^2} = \sigma^2$, and
    $\overline{X}_n - \mu = \frac{1}{n} \sum_{i} Y_i$. For $i \neq j$, $Y_i$ and $Y_j$ are independent,
    so $\expectation{Y_i Y_j} = \expectation{Y_i} \expectation{Y_j} = 0$, and expanding the square as
    in the proof of \cref{lem:pairwise-0903},
    \[\variance{\overline{X}_n}
        = \expectation{\Big(\frac{1}{n} \sum_{i=1}^{n} Y_i\Big)^2}
        = \frac{1}{n^2} \sum_{i=1}^{n} \expectation{Y_i^2}
        + \frac{1}{n^2} \sum_{i \neq j} \expectation{Y_i Y_j}
        = \frac{\sigma^2}{n}.\]
    Chebyshev's inequality (\cref{thm:chebyshev}) for $\overline{X}_n$ with $a = \eps$ gives
    $\prob{\abs{\overline{X}_n - \mu} \geq \eps} \leq \sfrac{\variance{\overline{X}_n}}{\eps^2}
        = \sfrac{\sigma^2}{(n \eps^2)}$, which tends to $0$ as $n \to \infty$.
\end{proof}
